% Figure: lift and drag coefficients of a spinning smooth cylinder in cross
% flow as functions of the spin ratio alpha = omega R / U (the ratio of the
% surface speed to the freestream). The lift coefficient rises steeply and
% roughly linearly at low spin, then rolls over toward a saturation plateau as
% the spin ratio grows; the drag coefficient falls from its stationary value at
% low spin before climbing again at high spin. Sketched schematically from the
% classical rotating-cylinder measurements; the curves are illustrative shapes,
% not a data fit.
\begin{figure}[htbp]
\centering
\begin{tikzpicture}
\begin{axis}[
  width=11cm, height=7cm,
  xlabel={spin ratio $\alpha = \omega R/U$},
  ylabel={coefficient},
  xmin=0, xmax=5, ymin=0, ymax=8,
  legend pos=north west,
  legend cell align=left,
  grid=both, grid style={enggray!18},
  minor grid style={enggray!8},
  tick label style={font=\scriptsize},
  label style={font=\small},
  legend style={font=\scriptsize},
]
% Lift coefficient: steep near-linear rise rolling over toward a plateau.
% Illustrative shape C_L = 7.0 * (1 - exp(-0.6 x)).
\addplot[engblue, very thick, domain=0:5, samples=100] {7.0*(1 - exp(-0.6*x))};
\addlegendentry{lift $C_L$}
% Drag coefficient: falls from the stationary value near 1.1, dips, then rises.
% Illustrative shape C_D = 1.1 - 0.7 x + 0.30 x^2, clamped positive.
\addplot[engred, very thick, domain=0:5, samples=100]
  {max(0.05, 1.1 - 0.7*x + 0.30*x^2)};
\addlegendentry{drag $C_D$}
% Mark the stationary-cylinder drag value on the left axis.
\node[font=\scriptsize, engred, anchor=west] at (axis cs:0.15,1.4)
  {$C_D(\alpha{=}0)\approx1.1$};
\end{axis}
\end{tikzpicture}
\caption[Magnus lift and drag of a spinning cylinder]{Lift and drag
coefficients of a smooth spinning cylinder in cross flow against the spin ratio
$\alpha = \omega R/U$, the surface speed in units of the freestream. Spin
displaces the two separation points asymmetrically, so the wake tilts and the
cylinder carries a transverse Magnus force; the lift climbs steeply and nearly
linearly at low spin before rolling toward a saturation plateau where further
spin drives little extra circulation. The drag first falls below its stationary
value as the moving surface delays separation on one side, then rises again at
high spin as the whole wake widens. Worked example~31 reads the low-spin part of
these curves for a rotor ship. The shapes are illustrative, not a data fit.}
\label{fig:vol03ch12:magnus-lift-drag}
\end{figure}
